\chapter{Supporting Constructions for Inheritance}
\label{s:inheritance}

This chapter gives two exact ways to determine what a parent must retain:
factorization of a boundary map and conditional prediction from a retained
history. Both distinguish an inherited response from a newly chosen parent
interaction.

\section{Linear boundary descent}

\propositionheading{Descent through a surjective retention}{sinh:linear-descent}
Let $E_j,E_i,Z_j,Z_i$ be finite-dimensional real vector spaces. Let
$C_j:E_j\to Z_j$ be surjective, let $C_i:E_i\to Z_i$ be linear, and let
$M_{ij}:E_j\to E_i$ be a linear boundary map. There is a linear map
$M_{IJ}^c:Z_j\to Z_i$ satisfying
\begin{equation}
 C_i\circ M_{ij}=M_{IJ}^c\circ C_j,
 \label{sinh:linear-square}
\end{equation}
if and only if
\begin{equation}
 \ker C_j\subseteq\ker(C_i\circ M_{ij}).
 \label{sinh:kernel-condition}
\end{equation}
When it exists, it is unique.

\begin{proof}
If \cref{sinh:linear-square} holds, every $v\in\ker C_j$ satisfies
$C_i(M_{ij}v)=0$. Conversely, define
$M_{IJ}^c(C_jv)=C_i(M_{ij}v)$. Two choices $v,w$ with $C_jv=C_jw$ differ
by an element of $\ker C_j$, which \cref{sinh:kernel-condition}
annihilates. The definition is therefore independent of the representative.
It is linear, and surjectivity of $C_j$ defines it everywhere. The same
surjectivity proves uniqueness.
\end{proof}

This is the factor theorem for linear maps. It covers deterministic linear
retentions only. For a stochastic sample retention the parent channel must be
a Markov kernel, and its existence is the garbling condition stated with the
channel square in the main text; a linear solution with negative entries does
not qualify.

For example, let $C_j(u,v)=u$ and $C_i$ be the identity on $\R$. The
map $M_{ij}(u,v)=u+v$ cannot descend: the discarded vector $(0,1)$ changes
the retained output. The map $M_{ij}(u,v)=2u$ descends to multiplication
by two. Agreement of two current input laws does not change either
conclusion, since the requirement concerns the boundary mechanism on its
declared domain.

The same criterion handles an additive-noise channel without requiring
Gaussian noise. Suppose the receiver content is $M_{ij}v+\varepsilon$,
where $\varepsilon$ has a fixed probability law $\nu$ independent of
$v$. If \cref{sinh:kernel-condition} holds, its retained law is the
distribution of
\begin{equation}
 M_{IJ}^c C_jv+C_i\varepsilon.
\end{equation}
Thus the parent inherits the pushed noise law $(C_i)_\#\nu$. A
state-dependent noise law requires its own factorization through $C_j$;
factoring the mean map alone is insufficient.

\section{Exact finite filtering}

Let $X_n$ be a Markov chain on a finite configuration space $\mathsf S$
with transition matrix $T$ and initial row distribution $p_0$. Let
$Z_n=R(X_n)$ for a map into a finite retained space $\mathsf S^c$.
For each $z$, let $P_z$ be the diagonal matrix whose entry at $x$ is
$1_{\{R(x)=z\}}$. These matrices are selectors of configurations, not
new dynamical transitions.

\propositionheading{The terminal conditional law retains every future prediction}{sinh:filter}
For a retained history $h_n=(z_0,\ldots,z_n)$ of positive probability,
define
\begin{equation}
 f_{h_n}=p_0P_{z_0}TP_{z_1}\cdots TP_{z_n},\qquad
 \widehat p_{h_n}=\frac{f_{h_n}}{f_{h_n}\mathbf 1}.
 \label{sinh:filter-law}
\end{equation}
Then $\widehat p_{h_n}$ is the conditional distribution of $X_n$ given
$h_n$, and
\begin{equation}
 \Pr(Z_{n+1}=z\mid h_n)=\widehat p_{h_n}TP_z\mathbf 1.
 \label{sinh:next-record}
\end{equation}
After a positive-probability next value $z$, the update is
\begin{equation}
 \widehat p_{(h_n,z)}
 =\frac{\widehat p_{h_n}TP_z}
        {\widehat p_{h_n}TP_z\mathbf 1}.
 \label{sinh:filter-update}
\end{equation}
For any finite future word, its probability is obtained by further
multiplication by the corresponding $TP_z$ matrices and by $\mathbf 1$.
The recursion is the forward algorithm for a function of a Markov chain
\citep{BaumPetrie1966,EphraimMerhav2002}. The normalized vector
$\widehat p_{h_n}$ is itself a Markov process on the simplex of laws on
$\mathsf S$, moving to \cref{sinh:filter-update} with probability
\cref{sinh:next-record}; it is the information state of a partially
observed Markov chain \citep{Astrom1965}. The proof uses the Markov
property at every propagation step, and time-dependent transition matrices
$T_n$ may replace $T$ without change.

\begin{proof}
The $x$ component of $p_0P_{z_0}$ is the joint probability of
$X_0=x$ and $Z_0=z_0$. Multiplication by $T$ propagates this joint mass
one step; multiplication by $P_{z_1}$ restricts it to the next retained
value. Induction gives the joint mass of $h_n$ and $X_n=x$ as the $x$
component of $f_{h_n}$. Its sum is the probability of $h_n$, so division
gives the conditional distribution. One further propagation and selection
gives \cref{sinh:next-record,sinh:filter-update}. Repeating
the operation gives every future word probability.
\end{proof}

A zero-probability history has zero unnormalized message and no conditional
law supplied by this formula. Because all matrices are nonnegative, every
extension of that history also has zero mass. The construction therefore
needs no arbitrary prediction on an impossible history.

\remarkheading{The filter beyond finite configuration spaces}{sinh:filter-general}
Let $\mathsf S$ be a measurable space, let the retained space $\mathsf Z$ be
countable, let $T_n$ be Markov kernels on $\mathsf S$, and let $R$ be
measurable. For a retained history $h$ of positive probability define the
measures $f_{(z_0)}(B)=p_0(B\cap R^{-1}\{z_0\})$ and
$f_{(h,z)}(B)=\int f_h(dx)T_n(x,B\cap R^{-1}\{z\})$. The argument of
\cref{sinh:filter} applies with measures in place of vectors: the normalized
measure $f_h/f_h(\mathsf S)$ is the conditional law of the current
configuration given $h$, the next retained value has probability
$f_{(h,z)}(\mathsf S)/f_h(\mathsf S)$, and the update holds at every history
of positive probability. When $\mathsf Z$ is uncountable, every history may
have probability zero. On standard Borel spaces the conditional laws then
exist as regular conditional distributions and the update is a jointly
measurable disintegration of the one-step joint law of next state and next
retained value \citep{Kallenberg2021}; the identities hold almost surely,
choices of versions on null sets cannot be avoided, and the remark above on
impossible histories has no analogue. The Markov property of the normalized
filter holds once a jointly measurable version of the update is fixed. On an
infinite configuration space the filter need not concentrate even when the
retained history determines the initial configuration: for the doubling map
$x\mapsto2x\bmod1$ with $R(x)=1_{\{x\geq1/2\}}$ and a uniform initial law,
the conditional law of $X_n$ given its retained history is uniform on an
interval of length $1/2$ for every $n$.

The filter is sufficient even when the retained process is not Markov.
For the deterministic cycle $1\to2\to3\to1$ with uniform initial law,
set $R(1)=R(2)=0$ and $R(3)=1$. The histories $(0,0)$ and $(1,0)$
both have positive probability and end at the same retained value.
The first ends at state $2$ and predicts $1$ next; the second ends at
state $1$ and predicts $0$. Hence the current retained value cannot
determine its next distribution. The terminal conditional law does: it is
the point mass $\delta_2$ after $(0,0)$ and $\delta_1$ after $(1,0)$. Here the
filter recovers the fine state, so it compresses nothing; a history-based
parent forgets information only when noise or nonobservability prevents the
history from determining the fine state.

\section{Successive retention of responses}

Suppose a declared joint law contains a fine state $X$ and a boundary
response $O$. Let $Z=R_{10}(X)$ and $W=R_{21}(Z)$, and assume all spaces
are standard Borel. For every bounded measurable response observable $g$,
conditional expectation gives
\begin{equation}
 \E[g(O)\mid W]
 =\E\bigl[\E[g(O)\mid Z]\mid W\bigr].
 \label{sinh:conditional-tower}
\end{equation}
To verify the identity, integrate either side over any event determined by
$W$. Such an event is also determined by $Z$, so both integrals equal
the integral of $g(O)$ over that event. This characterizes conditional
expectation up to null sets.

The identity concerns one declared joint law. Applying it to an equilibrium
law produces equilibrium-conditional responses. Applying it to the joint
law conditional on a retained history produces the corresponding predictive
responses. These are interchangeable only when their relevant conditional
laws agree. Exact boundary descent is stronger: it fixes the response on
every admitted fiber, independently of the conditional weights used there.
